\chapter{The Fourier Transform}\label{ch:fourier}

\begin{center}\itshape
A frequency integral can preserve length just as a change of orthonormal coordinates does.\\
On $L^2$, the Fourier transform loses no energy.
\end{center}

A string has discrete modes. A function on the whole line has a continuous range of frequencies, so a frequency sum becomes an integral. There is an immediate difficulty: the exponential waves themselves are not square-integrable, and the defining integral need not converge for an arbitrary $L^2$ input. Plancherel's theorem constructs the transform there by completion, preserving the inner product exactly.

Fourier's 1807 heat memoir and his 1822 \emph{Th\'eorie analytique de la chaleur} made trigonometric expansions central to analysis. The question was never only how to calculate coefficients. It was also what kind of convergence could recover the original function. Dirichlet, Riemann, Cantor and Lebesgue helped develop the distinctions needed to answer it. Plancherel's 1910 result gives the Hilbert-space answer on the line.

Volume~5, Chapter~8 treated the transform for continuous integrable functions and established its Gaussian identities, inversion under suitable hypotheses, Poisson summation and sampling. Here Lebesgue integration lets us separate the $L^1$ and $L^2$ theories, then connect them. Interpolation reaches the intermediate $L^p$ spaces. Measures, band-limited functions and lattices supply three settings in which the same transform answers different questions.

Fix an integer $d\ge1$ and use the normalisation
\[
\hat f(\xi)=\mathcal Ff(\xi)=\int_{\R^d}f(x)e^{-2\pi ix\cdot\xi}\dd x.
\]
An integral without a domain is over $\R^d$. Inner products use Lebesgue measure and are linear in the first variable.

\section{The transform on \texorpdfstring{$L^1$}{L1}}

\begin{proposition}\label{prop:ftl1}
For $f,g\in L^1(\R^d)$:
\begin{enumerate}
\item[(a)] $\hat f$ is continuous and $\|\hat f\|_\infty\le\|f\|_1$.
\item[(b)] $\hat f(\xi)\to0$ as $|\xi|\to\infty$ (Riemann--Lebesgue).
\item[(c)] Convolution becomes multiplication: $\widehat{f*g}=\hat f\hat g$.
\item[(d)] Translation $f(x-a)$ becomes $e^{-2\pi ia\cdot\xi}\hat f(\xi)$, and dilation $f(x/\lambda)$ becomes $\lambda^d\hat f(\lambda\xi)$ for $\lambda>0$.
\item[(e)] The multiplication formula is $\int\hat f\,g=\int f\,\hat g$.
\end{enumerate}
\end{proposition}
\begin{proof}
The integrand has absolute value $|f|$, giving the norm bound and, by dominated convergence, continuity in $\xi$. Changes of variables prove both rules in (d). Tonelli gives total absolute integral $\|f\|_1\|g\|_1$ for the double integrals in (c) and (e), so Fubini applies. For convolution put $u=x-y$; the exponential separates into its $u$ and $y$ factors, giving $\hat f(\xi)\hat g(\xi)$. Interchanging the $x$ and $\xi$ integrals gives (e).
For decay, take $\xi\ne0$ and $h=\xi/(2|\xi|^2)$. Translation contributes the phase $e^{-2\pi ih\cdot\xi}=-1$, so
\[
2\hat f(\xi)=\int(f(x)-f(x+h))e^{-2\pi ix\cdot\xi}\dd x.
\]
Its modulus is at most $\|f-f(\cdot+h)\|_1/2$. Since $|h|=1/(2|\xi|)\to0$, translation continuity in $L^1$ (Volume~7, Proposition~10.8) proves (b).
\end{proof}

The Gaussian $\gamma(x)=e^{-\pi|x|^2}$ is its own transform. This follows coordinate by coordinate from Volume~5, Theorem~8.3 and Fubini. Scaling and reversing the sign in the exponent give
\begin{equation}\label{eq:gaussft}
\int e^{-\pi\varepsilon^2|\xi|^2}e^{2\pi ix\cdot\xi}\dd\xi
=\varepsilon^{-d}e^{-\pi|x|^2/\varepsilon^2}=:G_\varepsilon(x).
\end{equation}
These non-negative kernels have integral one and concentrate at zero, so they form the approximate identity of Volume~7, Theorem~10.10.

\begin{theorem}[Fourier inversion]\label{thm:ftinv}
If $f,\hat f\in L^1(\R^d)$, then
\[
f(x)=\int\hat f(\xi)e^{2\pi ix\cdot\xi}\dd\xi
\]
for almost every $x$. If $f$ is continuous, the equality holds everywhere. In particular an $L^1$ function with zero transform vanishes almost everywhere.
\end{theorem}
\begin{proof}
Insert the Gaussian factor and call the integral $I_\varepsilon(x)$. Dominated convergence with bound $|\hat f|$ makes it tend, for every $x$, to the right side of the asserted formula. Expanding $\hat f$ and using Fubini is legitimate because $|f(y)|e^{-\pi\varepsilon^2|\xi|^2}$ has finite double integral. Equation~\eqref{eq:gaussft} then gives $I_\varepsilon=f*G_\varepsilon$.
The approximate-identity theorem gives $I_\varepsilon\to f$ in $L^1$. Along a sequence $\varepsilon_j\downarrow0$ it therefore converges almost everywhere, identifying the two limits. The inverse integral is continuous by dominated convergence. If $f$ is continuous too, their almost-everywhere equality is everywhere equality, since a nonzero continuous difference would stay nonzero on a positive-measure ball.
\end{proof}

\section{The Schwartz space}

Differentiation and multiplication exchange roles under the transform. To perform both repeatedly without losing integrability, use a space of smooth functions with rapid decay. For multi-indices $\alpha,\beta\in\N^d$, write $x^\alpha=\prod_jx_j^{\alpha_j}$ and $\partial^\beta=\prod_j\partial_j^{\beta_j}$. The \emph{Schwartz space} $\mathcal S(\R^d)$ consists of all smooth $f$ such that $\sup_x|x^\alpha\partial^\beta f(x)|<\infty$ for every pair of multi-indices. It contains compactly supported smooth functions and Gaussians and is closed under derivatives and polynomial multiplication. Bounding $(1+|x|^2)^df$ shows that $|f(x)|\le C(1+|x|^2)^{-d}$, hence $f\in L^p$ for every $1\le p\le\infty$.

\begin{proposition}\label{prop:schwartz}
The transform maps $\mathcal S(\R^d)$ bijectively to itself, with
\[
\widehat{\partial^\beta f}=(2\pi i\xi)^\beta\hat f,\qquad
\widehat{(-2\pi ix)^\alpha f}=\partial^\alpha\hat f.
\]
It satisfies $\mathcal F^2f(x)=f(-x)$ and $\mathcal F^4=I$ there.
\end{proposition}
\begin{proof}
Every polynomial multiple of $f$ is integrable, so differentiation under the integral (Volume~7, Theorem~8.17) gives the second formula and smoothness of $\hat f$. For a first partial derivative of $f$, use Fubini to integrate that coordinate first and integrate by parts on the line. Rapid decay removes its boundary terms and gives the factor $2\pi i\xi_j$. Repetition gives the first formula. Combining them expresses $\xi^\beta\partial^\alpha\hat f$ as a constant multiple of the transform of $\partial^\beta(x^\alpha f)$, which is integrable. The $L^1$ transform bound makes every such expression bounded, proving $\hat f\in\mathcal S$.
Inversion now applies everywhere to $f$ and gives $\mathcal F^2f(x)=f(-x)$. Applying it twice yields $\mathcal F^4=I$, so $\mathcal F^3$ is an inverse on this space.
\end{proof}

\section{Plancherel's theorem}

\begin{theorem}[Plancherel, 1910]\label{thm:plancherel}
There is a unique bounded operator $\mathcal F:L^2(\R^d)\to L^2(\R^d)$ agreeing with the integral transform on $L^1\cap L^2$. It is onto and preserves inner products, hence is unitary. Its inverse is $\mathcal F^{-1}g(x)=\mathcal Fg(-x)$, interpreted as $L^2$ classes. For every $f\in L^2$,
\[
\mathcal Ff=\lim_{R\to\infty}\int_{|x|\le R}f(x)e^{-2\pi ix\cdot\xi}\dd x
\quad\text{in }L^2.
\]
\end{theorem}
\begin{proof}
First take $f\in\mathcal S$ and $g=\overline{\hat f}\in\mathcal S$. Inversion gives $\hat g=\bar f$. The multiplication formula therefore gives
$\int|\hat f|^2=\int\hat f g=\int f\hat g=\int|f|^2$.
Smooth compactly supported functions are dense in $L^2$ (Volume~7, Theorem~10.11), so the isometry extends uniquely to its completion. Its range is complete and hence closed, and contains $\mathcal F(\mathcal S)=\mathcal S$, a dense subspace. Thus it is onto. Polarisation gives preservation of inner products. The proposed bounded inverse works on $\mathcal S$ and therefore everywhere by density.
To check agreement with the integral, let $f\in L^1\cap L^2$. First truncate in space, then mollify, choosing a sequence $f_j\in C_c^\infty$ tending to $f$ in both norms, using Volume~7's approximate-identity results. Its integral transforms converge uniformly to $\hat f$, while its $L^2$ transforms converge in $L^2$ to the extension. A subsequence converges almost everywhere, identifying the two limits. Finally $f\mathbf1_{|x|\le R}$ lies in $L^1\cap L^2$ by Cauchy--Schwarz on a finite-measure ball and tends to $f$ in $L^2$. Continuity gives the displayed truncation formula. Any bounded operator agreeing on $L^1\cap L^2$ agrees on a dense subspace, proving the stated uniqueness.
\end{proof}

\begin{beauty}
On the circle, a Fourier basis replaces a function by a square-summable sequence of coordinates. On the line, individual exponential waves do not belong to $L^2$, yet completion still gives a unitary transform of the whole space. We retain the inner-product identity without claiming an ordinary Hilbert basis made of non-square-integrable vectors. Chapter~\ref{ch:distributions} will give those waves a precise distributional role.
\end{beauty}

\section{The Hermite functions}

The Gaussian is fixed by $\mathcal F$. Other eigenfunctions arise by repeatedly applying a differential operator to it. On $\mathcal S(\R)$ define
\[
A=2\pi x+\frac{\dd}{\dd x},\qquad
A^\dagger=2\pi x-\frac{\dd}{\dd x}.
\]
Integration by parts makes them adjoint on this common domain. They are the annihilation and creation operators for the oscillator; this domain identity suffices here, without assuming an unbounded-operator spectral theorem.

\begin{theorem}[Hermite functions]\label{thm:hermite}
Set $g_n=(A^\dagger)^ne^{-\pi x^2}$ and $h_n=g_n/\|g_n\|_2$. Then $g_n=P_n(x)e^{-\pi x^2}$ for a polynomial of degree exactly $n$, and $(h_n)_{n\ge0}$ is a Hilbert basis of $L^2(\R)$. Its transform is
$\mathcal Fh_n=(-i)^nh_n$.
Thus the Fourier eigenvalues are $1,-i,-1,i$, each with infinite multiplicity.
\end{theorem}
\begin{proof}
Applying $A^\dagger$ to $Pe^{-\pi x^2}$ gives $(4\pi xP-P')e^{-\pi x^2}$. It raises degree by one, multiplying the leading coefficient by $4\pi$. This proves the polynomial shape and nonvanishing.
The Schwartz differentiation rules give $\mathcal F(A^\dagger f)=-iA^\dagger\hat f$. Since the Gaussian is fixed, induction gives the asserted eigenvalue for each $g_n$.
For orthogonality, direct differentiation gives $AA^\dagger-A^\dagger A=4\pi I$ and $Ag_0=0$. Induction then gives $Ag_n=4\pi n g_{n-1}$. The symmetric differential expression
$H=A^\dagger A+2\pi I=-\dd^2/\dd x^2+4\pi^2x^2$
satisfies $Hg_n=2\pi(2n+1)g_n$. Integration by parts on Schwartz functions makes eigenvectors at these distinct real values orthogonal. Normalising makes the family orthonormal.
To prove completeness, its first $n+1$ members span the same space as $x^ke^{-\pi x^2}$, $0\le k\le n$. If $f\in L^2$ is orthogonal to all of them, put
$\Phi(z)=\int f(x)e^{-\pi x^2}e^{-2\pi ixz}\dd x$.
Cauchy--Schwarz makes $\int|f(x)|e^{-\pi x^2}e^{2\pi|x||z|}\dd x$ finite. Thus the exponential series can be integrated term by term with absolute convergence, yielding
\[
\Phi(z)=\sum_{k\ge0}\frac{(-2\pi iz)^k}{k!}
\int f(x)x^ke^{-\pi x^2}\dd x=0.
\]
The integrable function $fe^{-\pi x^2}$ therefore has zero transform on the real axis. Inversion uniqueness makes it zero almost everywhere and hence $f=0$. This proves completeness. Finally, a bounded diagonal operator has spectrum equal to the closure of its diagonal values, by Exercise~\ref{ex:multiplication} for counting measure. Here those values are the four listed roots, with infinitely many indices in each class modulo four.
\end{proof}

\begin{figure}[ht]
\centering
\fig{ch05-hermite}
\caption{The Hermite functions $h_0,\dots,h_3$. Each is, up to the factor $(-i)^n$, its own Fourier transform; $h_n$ has $n$ zeros, and its energy $2\pi(2n+1)$ for the oscillator $H$ grows by equal
steps, the ``quanta'' of the harmonic oscillator.}
\label{fig:hermite}
\end{figure}

Products of these one-dimensional functions give the corresponding modes on $\R^d$. Fubini separates their transforms, multiplying their eigenvalues. Completeness follows by iterating the one-dimensional basis expansion; equivalently, finite products are dense by finite-measure rectangle approximation. Thus the same four spectral values occur in every dimension $d\ge1$.

\section{Interpolation: Riesz--Thorin and Hausdorff--Young}

The transform has norm at most one from $L^1$ to $L^\infty$ and norm one from $L^2$ to itself. Intermediate exponents are controlled by interpolation. Riesz's 1927 theorem and Thorin's 1938 complex-variable argument make this precise. The scalar estimate behind Thorin's proof is a maximum principle on a strip.

\begin{lemma}[three lines lemma]\label{lem:threelines}
Let $F$ be bounded and continuous on $0\le\operatorname{Re}z\le1$ and holomorphic in its interior. Suppose $|F(iy)|\le M_0$ and $|F(1+iy)|\le M_1$ for all real $y$, with $M_0,M_1>0$. Then
$|F(x+iy)|\le M_0^{1-x}M_1^x$ throughout the strip.
\end{lemma}
\begin{proof}
For $\varepsilon>0$ form
$G_\varepsilon(z)=F(z)M_0^{z-1}M_1^{-z}e^{\varepsilon z^2}$,
using real logarithms for the positive bases. Its modulus is at most $e^\varepsilon$ on the two vertical edges. Boundedness of $F$ and the factor $e^{\varepsilon(x^2-y^2)}$ make it tend uniformly to zero as $|y|\to\infty$, for $0\le x\le1$. Choose a rectangle tall enough that its horizontal edges also have modulus at most $e^\varepsilon$. The maximum principle (Volume~5, Theorem~9.19) gives this bound inside. For any fixed strip point such a rectangle can be chosen, giving
$|F(z)|M_0^{x-1}M_1^{-x}\le e^\varepsilon e^{-\varepsilon(x^2-y^2)}$.
Let $\varepsilon\downarrow0$ to obtain the claim.
\end{proof}

\begin{theorem}[Riesz--Thorin]\label{thm:rieszthorin}
Let $(X,\mu),(Y,\nu)$ be sigma-finite. Let $T$ be linear on the simple functions with finite-measure support on $X$, taking values in measurable functions on $Y$. Suppose
$\|Tf\|_{q_j}\le M_j\|f\|_{p_j}$ for $j=0,1$, where $1\le p_j,q_j\le\infty$.
For $0<\theta<1$ set
\[
\frac1{p_\theta}=\frac{1-\theta}{p_0}+\frac\theta{p_1},\qquad
\frac1{q_\theta}=\frac{1-\theta}{q_0}+\frac\theta{q_1}.
\]
If $p_\theta<\infty$ and $1<q_\theta<\infty$, then
$\|Tf\|_{q_\theta}\le M_0^{1-\theta}M_1^\theta\|f\|_{p_\theta}$.
The map extends uniquely to $L^{p_\theta}\to L^{q_\theta}$ with that norm bound.
\end{theorem}
\begin{proof}
Write $p=p_\theta$, $q=q_\theta$ and let $q'$ be conjugate to $q$. If either $M_j=0$, $T$ vanishes on its whole simple-function domain and the conclusion follows by density. Assume both are positive.

\emph{Dual tests.} For $h\in L^q$,
$\|h\|_q=\sup|\int hg|$ over simple finite-support $g$ with $\|g\|_{q'}\le1$. H\"older gives the upper bound. For nonzero $h$, the function $g_0=|h|^{q-1}\overline{\operatorname{sgn}h}/\|h\|_q^{q-1}$ has norm one and gives equality. Approximate it in $L^{q'}$ by simple finite-support functions and rescale their norms to at most one. H\"older passes the pairings to the limit. The zero case is immediate. Here $q'<\infty$, as required for density.

\emph{An analytic finite sum.} Take simple $f=\sum_ja_j\mathbf1_{A_j}$ and $g=\sum_kb_k\mathbf1_{B_k}$, each with disjoint finite-measure level sets, and normalise $\|f\|_p=\|g\|_{q'}=1$. Omit zero coefficients. Define the affine reciprocal exponents
\[
\frac1{p(z)}=\frac{1-z}{p_0}+\frac z{p_1},\qquad
\frac1{q'(z)}=\frac{1-z}{q'_0}+\frac z{q'_1},
\]
and put
\[
f_z=\sum_j|a_j|^{p/p(z)}\operatorname{sgn}(a_j)\mathbf1_{A_j},\qquad
g_z=\sum_k|b_k|^{q'/q'(z)}\operatorname{sgn}(b_k)\mathbf1_{B_k}.
\]
The pairing $F(z)=\int(Tf_z)g_z\dd\nu$ is a finite linear combination of entire exponentials of affine functions. Its coefficients $\int T\mathbf1_{A_j}\mathbf1_{B_k}$ are finite by an endpoint H\"older bound and the finite measure of $B_k$. Thus $F$ is entire and bounded on the closed strip, since the moduli of those exponentials depend only on $\operatorname{Re}z$ there.
At $z=\theta$, $f_z=f$ and $g_z=g$. At $z=iy$, the moduli are $|f|^{p/p_0}$ and $|g|^{q'/q_0'}$, giving endpoint norms one. If an endpoint exponent is infinite, the corresponding modulus is one on its support, so the same norm statement holds. H\"older and the assumed operator bound give $|F(iy)|\le M_0$; the other edge similarly gives $M_1$.

The three-lines lemma now bounds $|\int(Tf)g|$ by $M_0^{1-\theta}M_1^\theta$. For this simple $f$, the endpoint assumptions put $Tf$ in $L^{q_0}\cap L^{q_1}\subseteq L^q$, by the interpolation inequality for one function (Volume~7, Exercise~10.3). The dual tests therefore give the asserted norm bound. Homogeneity removes the normalisation. Simple finite-support functions are dense in $L^p$ because $p<\infty$ (Volume~7, Proposition~10.6), and the target is complete, giving the unique extension.
\end{proof}

\begin{corollary}[Hausdorff--Young]\label{cor:hausdorffyoung}
For $1\le p\le2$ and $1/p+1/p'=1$, every $f\in L^1\cap L^p(\R^d)$ satisfies $\|\hat f\|_{p'}\le\|f\|_p$. The Fourier transform extends boundedly from $L^p$ to $L^{p'}$ with norm at most one.
\end{corollary}
\begin{proof}
At $p=1$ use the $L^1$ bound, and at $p=2$ use Plancherel. Otherwise interpolate between $(p_0,q_0)=(1,\infty)$ and $(p_1,q_1)=(2,2)$, both with norm one, taking $\theta=2-2/p$. Then $p_\theta=p$ and $q_\theta=p'$. To check consistency for $f\in L^1\cap L^p$, choose the same simple finite-support approximants in both norms, by truncating the support and the values and approximating the bounded remainder. Integral transforms converge uniformly, while the extension converges in $L^{p'}$, hence along a subsequence almost everywhere. The limits agree. Density gives the extension on the whole $L^p$ space.
\end{proof}

\begin{figure}[ht]
\centering
\fig{ch05-riesz}
\caption{The Riesz diagram: an operator bounded from $L^p$ to $L^q$ is a point $(\frac1p,\frac1q)$ of the square. The set of such points is convex (Theorem~\ref{thm:rieszthorin}). For $\mathcal F$, the
endpoints $(1,0)$ (Proposition~\ref{prop:ftl1}) and $(\frac12,\frac12)$ (Plancherel) give the Hausdorff--Young segment $\frac1q=1-\frac1p$.}
\label{fig:riesz}
\end{figure}

For $1<p<2$ the best Fourier operator norm is smaller than one. Beckner's 1975 theorem gives the exact value
$\bigl(p^{1/p}/p'^{1/p'}\bigr)^{d/2}$,
attained by Gaussians. This sharp theorem is quoted, not proved by the interpolation argument. Exercise~\ref{ex:hysharp} computes the Gaussian ratio and distinguishes that calculation from the missing global upper bound. Riesz--Thorin also yields Young's convolution inequality by interpolating the endpoint convolution estimates (Volume~7, Proposition~10.9 and Exercise~10.1).

\section{Measures and characteristic functions}

For a finite Borel measure $\mu$ on $\R^d$, define $\hat\mu(\xi)=\int e^{-2\pi ix\cdot\xi}\dd\mu(x)$. It is bounded and continuous by dominated convergence. A random vector with this law has characteristic function
$\varphi_X(t)=\mathbb E e^{it\cdot X}=\hat\mu(-t/(2\pi))$,
matching Volume~4, Chapter~6 after the change of normalisation.

\begin{theorem}\label{thm:measuniq}
Two finite Borel measures on $\R^d$ with identical Fourier transforms are equal. In particular a characteristic function determines its probability law.
\end{theorem}
\begin{proof}
With a Schwartz input $g$, Fubini rewrites the pairing as $\int\hat g\dd\mu=\int g(\xi)\hat\mu(\xi)\dd\xi$. Equality of the transforms makes these pairings equal for $\mu$ and $\nu$. Since $\mathcal F$ maps $\mathcal S$ onto itself, the measures agree on every Schwartz test function, in particular every $C_c^\infty$ test.
If $h\in C_c$, mollification gives uniformly convergent $C_c^\infty$ approximants, with supports in one compact set (Volume~7, Theorem~10.10(b)). Finiteness of both measures passes equality to $h$. For a bounded open box $U$, take continuous compactly supported $0\le h_j\uparrow\mathbf1_U$, for example truncated multiples of the distance to $U^c$. Monotone convergence gives equality of the two box measures. Their total masses agree at transform value zero. Bounded open boxes, together with the empty set, form a generating $\pi$-system, so finite-measure uniqueness (Volume~7, Theorem~7.5) gives equality on all Borel sets.
\end{proof}

\begin{proposition}\label{prop:cauchylaw}
For $a>0$ on the line, $e^{-2\pi a|x|}$ has transform $p_a(\xi)=a/(\pi(a^2+\xi^2))$, and $\hat p_a(\xi)=e^{-2\pi a|\xi|}$. The Cauchy probability density $p_a$ therefore has characteristic function $e^{-a|t|}$.
\end{proposition}
\begin{proof}
Integrate separately on the two half-lines to obtain
\[
\frac1{2\pi(a+i\xi)}+\frac1{2\pi(a-i\xi)}
=\frac a{\pi(a^2+\xi^2)}.
\]
Both functions are integrable and continuous, and $p_a$ is even. Inversion therefore gives its transform as the original exponential. Substitution $\xi=-t/(2\pi)$ gives the characteristic function. Also $\int p_a=\hat p_a(0)=1$, confirming its normalisation.
\end{proof}

Multiplying these transforms gives $\hat p_s\hat p_t=\hat p_{s+t}$. Convolution and uniqueness then prove $p_s*p_t=p_{s+t}$, the identity from Volume~7, Exercise~10.2. It also shows why averages of Cauchy observations do not sharpen their distribution.

\begin{corollary}\label{cor:cauchymean}
For independent Cauchy variables $X_1,\ldots,X_n$ of density $p_a$, their mean $n^{-1}\sum_jX_j$ has the same density $p_a$.
\end{corollary}
\begin{proof}
Independence multiplies characteristic functions (Volume~4, Chapter~6), giving $(e^{-a|t|/n})^n=e^{-a|t|}$. The uniqueness theorem identifies the law.
\end{proof}

The absolute first moment of this density diverges, so the usual integrable-variable law of large numbers does not apply. The averaging identity makes the absence of concentration exact. Volume~11 uses characteristic functions to study convergence in law, including L\'evy's continuity theorem and the central limit theorem; these are later applications, not inputs here.

\section{Band-limited functions}

A square-integrable function is \emph{band-limited} to $W>0$ when its transform vanishes almost everywhere outside $[-W,W]$. Define
\[
PW_W=\{f\in L^2(\R):\mathcal Ff=0\text{ a.e. outside }[-W,W]\}.
\]
It is closed, since its image under the unitary transform is the closed subspace $L^2([-W,W])$. The integral representative in the next theorem gives meaningful point samples of these $L^2$ classes.

\begin{theorem}[sampling in Hilbert space]\label{thm:sampling}
Every $f\in PW_W$ has a unique continuous representative, extending to the entire function
$f(z)=\int_{-W}^W\hat f(\xi)e^{2\pi iz\xi}\dd\xi$.
The functions
\[
\varphi_n(x)=\sqrt{2W}\operatorname{sinc}(2Wx-n),\quad n\in\Z,
\qquad\operatorname{sinc}t=\frac{\sin\pi t}{\pi t},\quad\operatorname{sinc}0=1,
\]
form a Hilbert basis of $PW_W$. For its representative,
\[
f(x)=\sum_{n\in\Z}f\left(\frac n{2W}\right)\operatorname{sinc}(2Wx-n),
\]
with $L^2$ and uniform convergence on $\R$.
\end{theorem}
\begin{proof}
Finite frequency support and Cauchy--Schwarz put $\hat f$ in $L^1$. The integral defines an entire function: the exponential series converges absolutely under the integral, uniformly for $z$ in compact sets, since $|\xi|\le W$. Plancherel's consistency identifies its real restriction with the given $L^2$ class. Two continuous representatives agreeing almost everywhere agree everywhere.
The functions $\psi_n(\xi)=(2W)^{-1/2}e^{-2\pi in\xi/(2W)}\mathbf1_{[-W,W]}(\xi)$ are a Hilbert basis on that interval, by rescaling Theorem~\ref{thm:fourierl2}. Direct integration gives $\mathcal F^{-1}\psi_n=\varphi_n$, so unitarity transports the basis to $PW_W$. Its coefficients are
\[
\langle f,\varphi_n\rangle=\langle\hat f,\psi_n\rangle
=\frac1{\sqrt{2W}}f\left(\frac n{2W}\right),
\]
giving the sampling expansion in $L^2$. For any $g\in PW_W$, its representative obeys
$|g(x)|\le\|\hat g\|_1\le\sqrt{2W}\|g\|_2$ uniformly in real $x$. Apply this to the basis-expansion remainder. Its $L^2$ norm tends to zero, so its uniform norm does too.
\end{proof}

The integrable sampling theorem in Volume~5, Theorem~8.15 becomes an orthonormal-coordinate identity here. The samples themselves are scalar multiples of basis coefficients.

\begin{corollary}\label{cor:nocompact}
If $f\in L^2(\R)$ and both $f$ and $\hat f$ vanish almost everywhere outside bounded sets, then $f=0$ almost everywhere.
\end{corollary}
\begin{proof}
The frequency bound gives the entire representative from the sampling theorem. The spatial bound makes it zero almost everywhere on a real open ray. Continuity makes it zero everywhere on that ray. The complex identity theorem (Volume~5, Theorem~9.18) now makes the entire function identically zero.
\end{proof}

\begin{think}
Exact finite duration and exact finite bandwidth cannot occur together for a nonzero signal. In practice a small tail can be ignored. How could eigenvalues of a time restriction followed by a frequency restriction quantify that approximation? Compare the compact min--max principle with this exact impossibility statement.
\end{think}
\begin{think}
A camera taking $24$ frames per second sees $23$ equally spaced spoke passages per second as a slow backward motion. Samples of frequencies $23$ and $-1$ agree at times $m/24$. Which bandwidth assumption in the sampling theorem would distinguish them?
\end{think}

\section{Poisson summation on lattices}

Write a full lattice as $\Lambda=A\Z^d$ with invertible real $A$. Its covolume $|\det A|$ is the volume of a fundamental parallelepiped. A vector pairs integrally with every lattice vector exactly when $A^{\mathsf T}w\in\Z^d$, so the dual lattice is $\Lambda^*=A^{-T}\Z^d$.

The spatial sum below converges uniformly on compact sets. Global uniform convergence of its ordinary series would be false: for a Gaussian and $\Lambda=\Z$, a tail containing index $m$ has value at least one at $x=-m$, however large that index is. The frequency-side series does converge uniformly over all $x$. Their sums still define the same continuous periodic function.

\begin{theorem}[Poisson summation]\label{thm:poissonlattice}
For $f\in\mathcal S(\R^d)$ and every $x\in\R^d$,
\[
\sum_{v\in\Lambda}f(x+v)
=\frac1{\operatorname{covol}\Lambda}
\sum_{w\in\Lambda^*}\hat f(w)e^{2\pi iw\cdot x}.
\]
The left series converges absolutely and locally uniformly in $x$; the right converges absolutely and uniformly on all $\R^d$. In particular their values at $x=0$ give the lattice summation identity.
\end{theorem}
\begin{proof}
First take $\Lambda=\Z^d$. On any fixed compact set of $x$, Schwartz decay bounds $|f(x+m)|$ by a constant times $(1+|m|)^{-d-1}$. This summable bound gives local uniform absolute convergence. The sum $F$ is continuous and periodic. Integrate over a unit cube and unfold the absolutely integrable sum to obtain its Fourier coefficients $\hat f(k)$:
\[
\int_{[0,1)^d}F(x)e^{-2\pi ik\cdot x}\dd x
=\int_{\R^d}f(x)e^{-2\pi ik\cdot x}\dd x.
\]
Schwartz decay of $\hat f$ gives $\sum_k|\hat f(k)|<\infty$, so its frequency series converges absolutely uniformly on the whole space. It has the same Fourier coefficients as $F$. Completeness of the torus exponentials, proved in Chapter~\ref{ch:hilbert} before Theorem~\ref{thm:kroneckerweyl}, identifies the two functions in $L^2$ on a cube; continuity and periodicity give equality everywhere.
For general $A$, put $g(y)=f(Ay)$, still Schwartz. Change variables to obtain $\hat g(\eta)=|\det A|^{-1}\hat f(A^{-T}\eta)$. Apply the integer-lattice formula at $y=A^{-1}x$. The indices become $v=Am$ and $w=A^{-T}k$, and $k\cdot A^{-1}x=w\cdot x$, giving the asserted formula and both convergence statements.
\end{proof}

\begin{corollary}[theta of a lattice]\label{cor:thetalattice}
For $t>0$ put $\theta_\Lambda(t)=\sum_{v\in\Lambda}e^{-\pi t|v|^2}$. Then
\[
\theta_\Lambda(t)=\frac{t^{-d/2}}{\operatorname{covol}\Lambda}
\theta_{\Lambda^*}(1/t),
\qquad
\theta_\Lambda(t)\sim\frac{t^{-d/2}}{\operatorname{covol}\Lambda}\quad(t\downarrow0).
\]
\end{corollary}
\begin{proof}
Apply Poisson summation at zero to the scaled Gaussian, whose transform is $t^{-d/2}e^{-\pi|\xi|^2/t}$. For $0<t\le1$ the dual summands are bounded by the summable $e^{-\pi|w|^2}$ and tend to zero except at $w=0$. Dominated convergence for this counting measure gives $\theta_{\Lambda^*}(1/t)\to1$, proving the asymptotic.
\end{proof}

For $\Z$, the identity is Volume~5's theta transformation (Corollary~8.10), used there in the zeta functional equation. The leading term records the lattice density $1/\operatorname{covol}\Lambda$ through Gaussian weights. For the self-dual $E_8$ lattice it reads $\theta(t)=t^{-4}\theta(1/t)$. Volume~13 studies the modular-form setting of these identities; Viazovska's sphere-packing proof uses further structure beyond this transformation alone.

\section{Differentiating beyond functions}

Completion preserved Fourier energy for all $L^2$ inputs. The same map recovered probability laws from characteristic functions and point values from a bandwidth constraint. On lattices, periodising a rapidly decreasing function produced a frequency sum on the dual lattice, with the convergence distinctions made explicit.

Differentiation still asks for more. A discontinuous function need not have an ordinary derivative, but integration against smooth test functions allows a derivative to be defined by integration by parts. Chapter~\ref{ch:distributions} builds distributions and extends the Fourier transform to them.

\begin{figure}[ht]
\centering\fig{ch10-rw-samples}
\caption{At bandwidth $W=1/2$, the sinc function is one sampling-basis vector. Its integer samples vanish except for the value one at zero; its long spatial tails are compatible with a compact frequency interval.}
\label{fig:v8c10samples}
\end{figure}
\begin{figure}[ht]
\centering\fig{ch10-rw-gaussian}
\caption{The functions $f(x)=e^{-4\pi x^2}$ and $\hat f(\xi)=\tfrac12e^{-\pi\xi^2/4}$ show reciprocal Gaussian widths with the required amplitude change. A narrow input has a wider transform.}
\label{fig:v8c10gaussian}
\end{figure}

\section{Exercises}

\begin{exercise}[$\star$]\label{ex:cothsum}
Use Volume~5, Theorem~8.9 on $e^{-2\pi a|x|}$ with $a>0$ to obtain $\sum_{n\in\Z}(n^2+a^2)^{-1}=(\pi/a)\coth(\pi a)$. Recover the Basel sum $\sum_{n\ge1}n^{-2}=\pi^2/6$ by letting $a$ decrease to zero.
\end{exercise}
\begin{exercise}[$\star$]\label{ex:heatft}
\emph{Heat on the whole space.} For Schwartz $f$ and $t>0$, define
$u(t,x)=\int e^{-4\pi^2t|\xi|^2}\hat f(\xi)e^{2\pi ix\cdot\xi}\dd\xi$.
Prove $u(t,\cdot)=f*K_t$ with $K_t(x)=(4\pi t)^{-d/2}e^{-|x|^2/(4t)}$. Verify $\partial_tu=\Delta u$ and uniform convergence to $f$ as $t\downarrow0$. Show that its $L^2$ norm is non-increasing in $t$ and its supremum norm is at most $(4\pi t)^{-d/2}\|f\|_1$.
\end{exercise}
\begin{exercise}[$\star\star$]\label{ex:bochner}
\emph{Positive definite functions.} Call $\varphi:\R^d\to\C$ positive definite when $\sum_{i,j}c_i\bar c_j\varphi(\xi_i-\xi_j)\ge0$ for every finite choice of points and coefficients. (a) Prove this property for transforms of finite positive measures. (b) If $\varphi$ is also continuous and integrable, prove that its inverse Fourier integral is non-negative at every point. (c) Prove $|\varphi(\xi)|\le\varphi(0)$; use it to exclude $\xi^2e^{-\pi\xi^2}$ on the line, and prove that $\max(1-|\xi|,0)$ is positive definite. The full converse, Bochner's theorem for continuous positive definite functions, is not proved here.
\end{exercise}
\begin{exercise}[$\star\star$]\label{ex:paleywiener}
\emph{One direction of Paley--Wiener.} If $f\in L^2(\R)$ is supported almost everywhere in $[-R,R]$ with $R>0$, show that its transform has an entire extension bounded by $\sqrt{2R}\|f\|_2e^{2\pi R|\operatorname{Im}\zeta|}$. Compute the extension for the interval indicator and prove that no smaller exponential rate works uniformly for this example. The converse requires additional contour arguments, of the kind developed in Chapter~\ref{ch:complex}.
\end{exercise}
\begin{exercise}[$\star$]\label{ex:hermitecalc}
\emph{Hermite calculations.} Derive $\|g_n\|_2^2=4\pi n\|g_{n-1}\|_2^2=(4\pi)^nn!/\sqrt2$ and compute $g_1,g_2$. Verify $\mathcal Fg_1=-ig_1$ from the Schwartz formulas. Prove $h_n(-x)=(-1)^nh_n(x)$. For real $\alpha$, define $\mathcal F_\alpha f=\sum_ne^{-in\alpha}\langle f,h_n\rangle h_n$ and prove the unitary group law, with $\mathcal F_{\pi/2}=\mathcal F$ and $\mathcal F_\pi f=f(-\cdot)$. This is the fractional Fourier transform.
\end{exercise}
\begin{exercise}[$\star$]\label{ex:hysharp}
For $\gamma(x)=e^{-\pi x^2}$ and $1<p<2$, calculate $\|\gamma\|_p$ and $\|\hat\gamma\|_{p'}$ and prove their ratio is strictly less than one. Does this one-input calculation alone show that the operator-norm constant one can be improved for every input? Identify the additional assertion supplied by Beckner's sharp theorem quoted after Corollary~\ref{cor:hausdorffyoung}.
\end{exercise}
\begin{exercise}[$\star\star$]\label{ex:hexlattice}
\emph{Hexagonal duality.} Let $\Lambda$ have basis $(1,0),(1/2,\sqrt3/2)$. Compute its covolume and a dual basis. Prove that $\Lambda^*$ is $\Lambda$ rotated by $90$ degrees and scaled by $2/\sqrt3$. Derive $\theta_\Lambda(t)=2(\sqrt3t)^{-1}\theta_\Lambda(4/(3t))$. Examine the fixed point $t=2/\sqrt3$.
\end{exercise}

\begin{exercise}[$\star\star$]\label{ex:v8c10a}
For $f_a(x)=e^{-\pi ax^2}$ with $a>0$, calculate its transform, its $L^2$ norm and the two norms $\|xf_a\|_2$, $\|\xi\hat f_a\|_2$. Verify $\|xf_a\|_2\|\xi\hat f_a\|_2=\|f_a\|_2^2/(4\pi)$.
\end{exercise}
\begin{exercise}[$\star$]\label{ex:v8c10b}
In $PW_{1/2}$ take $f(x)=\operatorname{sinc}x+2\operatorname{sinc}(x-1)$. Find its integer samples, Fourier transform, $L^2$ norm and value at $x=1/2$. Derive the general sample-energy identity $\|g\|_2^2=(2W)^{-1}\sum_{n\in\Z}|g(n/(2W))|^2$ for $g\in PW_W$.
\end{exercise}
\begin{exercise}[$\star\star$]\label{ex:v8c10c}
Let $\mu=(\delta_{-1}+\delta_1)/2$. Compute $\hat\mu$ and every convolution power $\mu^{*n}$. Explain why its transform does not tend to zero at infinity, and why this does not contradict Riemann--Lebesgue for $L^1$ functions.
\end{exercise}
\begin{exercise}[$\star\star$]\label{ex:v8c10d}
Let $\Lambda=2\Z\times3\Z$. Compute its dual and covolume. With $\theta(s)=\sum_{n\in\Z}e^{-\pi sn^2}$, prove
\[
\theta(4t)\theta(9t)=\frac1{6t}\theta(1/(4t))\theta(1/(9t)).
\]
Find the leading behaviour as $t\downarrow0$, and show that $\theta_\Lambda(t)=(6t)^{-1}\theta_\Lambda(1/(36t))$ has fixed point $t=1/6$.
\end{exercise}

\section{Hints}

\begin{hint}{ex:cothsum}
The exponential's transform is $p_a$. Check the decay hypotheses of Volume~5's theorem, then sum the spatial geometric series.
\end{hint}
\begin{hint}{ex:heatft}
Use the Gaussian identity with $\varepsilon^2=4\pi t$. The transformed time derivative and Laplacian both multiply by $-4\pi^2|\xi|^2$.
\end{hint}
\begin{hint}{ex:bochner}
Establish the two-point bound before using it in part (b). Integrate the positive kernel against a Gaussian times $e^{2\pi ix\cdot\xi}$, then complete the square.
\end{hint}
\begin{hint}{ex:paleywiener}
Use Cauchy--Schwarz on the finite support. Along the imaginary axis the interval indicator produces a hyperbolic sine.
\end{hint}
\begin{hint}{ex:hermitecalc}
Move $A^\dagger$ across the inner product and use $Ag_n=4\pi ng_{n-1}$. Even and odd parity alternate under $A^\dagger$.
\end{hint}
\begin{hint}{ex:hysharp}
Put $a=1/p$ and compare $a\log a$ with $(1-a)\log(1-a)$ by strict concavity. An operator norm is a supremum over all inputs.
\end{hint}
\begin{hint}{ex:hexlattice}
Use the inverse transpose of the basis matrix. Differentiating the theta identity at its fixed point yields more than merely substituting there.
\end{hint}
\begin{hint}{ex:v8c10a}
Differentiate the Gaussian integral with respect to $a$ to obtain a second moment, then apply the scaling formula to the transform.
\end{hint}
\begin{hint}{ex:v8c10b}
At this bandwidth the shifted sinc functions already have norm one. Their samples are Kronecker deltas.
\end{hint}
\begin{hint}{ex:v8c10c}
A sum of $n$ independent signs takes value $n-2k$ with binomial multiplicity. Test the transform at integer frequencies.
\end{hint}
\begin{hint}{ex:v8c10d}
Separate the two coordinate sums. The dual rectangle is a rotated copy of the original scaled by $1/6$.
\end{hint}

\section{Solutions}

\begin{solution}{ex:cothsum}
\idea Choose a function whose samples form a geometric series while its transform gives the desired rational summands.
\writeup Both the continuous exponential and its transform $p_a$ are bounded by constants times $(1+|x|)^{-2}$. Thus Volume~5, Theorem~8.9 applies, giving
\[
\sum_{m\in\Z}\frac a{\pi(a^2+m^2)}
=1+2\sum_{n\ge1}e^{-2\pi an}
=\frac{1+e^{-2\pi a}}{1-e^{-2\pi a}}=\coth(\pi a).
\]
Multiply by $\pi/a$. After removing the $m=0$ term and dividing by two, the right side is
$\tfrac12((\pi/a)\coth(\pi a)-a^{-2})$.
The expansion $\coth z=z^{-1}+z/3+O(z^3)$, obtained by dividing the Taylor series for cosh and sinh, makes its limit $\pi^2/6$. Monotone convergence on the sum $\sum_{n\ge1}(n^2+a^2)^{-1}$ as $a\downarrow0$ gives Basel's identity.
\end{solution}

\begin{solution}{ex:heatft}
\idea Each frequency evolves by a scalar decaying exponential, which makes both the equation and the energy estimate explicit.
\writeup The Gaussian formula gives $\hat K_t=e^{-4\pi^2t|\xi|^2}$. Convolution followed by inversion therefore identifies $f*K_t$ with the specified integral everywhere, since the transformed product is integrable and the convolution is continuous.
Differentiating in time introduces $-4\pi^2|\xi|^2$, as does the spatial Laplacian. All required derivatives are justified locally for $t>0$ by integrable polynomial multiples of the Schwartz function $\hat f$. Thus $\partial_tu=\Delta u$. Inversion for $f$ gives the uniform error bound
$|u(t,x)-f(x)|\le\int(1-e^{-4\pi^2t|\xi|^2})|\hat f(\xi)|\dd\xi\to0$
by dominated convergence.
Plancherel gives $\|u(t,\cdot)\|_2^2=\int e^{-8\pi^2t|\xi|^2}|\hat f(\xi)|^2\dd\xi$, non-increasing in $t$. Finally the convolution estimate is $|u(t,x)|\le\|f\|_1\|K_t\|_\infty=(4\pi t)^{-d/2}\|f\|_1$.
\end{solution}

\begin{solution}{ex:bochner}
\idea Finite positive quadratic sums remain positive after integration; Gaussian testing then isolates an inverse transform value.
\writeup (a) Expansion gives
\[
\sum_{i,j}c_i\bar c_j\hat\mu(\xi_i-\xi_j)
=\int\left|\sum_ic_ie^{-2\pi ix\cdot\xi_i}\right|^2\dd\mu(x)\ge0.
\]
Before (b), establish the bound needed from (c). The two-point quadratic condition forces $\varphi(-\xi)=\overline{\varphi(\xi)}$, $\varphi(0)\ge0$ and non-negative determinant of the $2\times2$ matrix with these entries. Hence $|\varphi(\xi)|\le\varphi(0)$ everywhere.
(b) For continuous compactly supported $g$, the quadratic integral
$Q(g)=\int\!\int\varphi(\xi-\eta)g(\xi)\overline{g(\eta)}\dd\xi\dd\eta$
is a limit of non-negative finite Riemann sums. Therefore it is real and non-negative. The bound just proved allows continuous cutoffs and dominated convergence to extend this conclusion to
$g(\xi)=e^{2\pi ix\cdot\xi}e^{-\pi\varepsilon^2|\xi|^2}$.
Put $u=\xi-\eta$ and use $|\eta+u|^2+|\eta|^2=2|\eta+u/2|^2+|u|^2/2$. Integration over $\eta$ gives
\[
0\le Q(g)=(2\varepsilon^2)^{-d/2}
\int\varphi(u)e^{2\pi ix\cdot u}e^{-\pi\varepsilon^2|u|^2/2}\dd u.
\]
Multiply by the positive prefactor's inverse and let $\varepsilon\downarrow0$. Integrability of $\varphi$ supplies domination, proving that the inverse Fourier integral is non-negative.
(c) The bound excludes $\xi^2e^{-\pi\xi^2}$, which vanishes at zero but not everywhere. For the triangle $\Delta$, write $\Delta=\mathbf1_{[-1/2,1/2]}*\mathbf1_{[-1/2,1/2]}$. Its transform is $\operatorname{sinc}^2$, an integrable non-negative function. Inversion and evenness give $\Delta=\widehat{\operatorname{sinc}^2}$. It is thus the transform of a finite positive density, so (a) proves positive definiteness.
Bochner's full theorem says every continuous positive definite function is the transform of a finite positive measure. Parts (a)--(c) establish the stated special consequences, without proving that full converse.
\end{solution}

\begin{solution}{ex:paleywiener}
\idea A bounded spatial support limits the growth of the complex exponential in the transform integral.
\writeup For $|x|\le R$, the exponential modulus is $e^{2\pi x\operatorname{Im}\zeta}\le e^{2\pi R|\operatorname{Im}\zeta|}$. Cauchy--Schwarz gives $\int|f|\le\sqrt{2R}\|f\|_2$, yielding the stated bound. Expanding the exponential under that finite-support integral, with compact-uniform domination in $\zeta$, proves the extension is entire. On the real axis it is the integral transform, agreeing with Plancherel.
For the indicator, direct integration gives $F(\zeta)=\sin(2\pi R\zeta)/(\pi\zeta)$, with the removable value $2R$ at zero. Along $\zeta=iy$, its positive value is $\sinh(2\pi Ry)/(\pi y)$, asymptotic to $e^{2\pi R|y|}/(2\pi|y|)$. Any proposed bound $Ce^{c|y|}$ with $c<2\pi R$ would fail as $|y|\to\infty$, proving optimality of that exponential rate for the example.
\end{solution}

\begin{solution}{ex:hermitecalc}
\idea The commutator supplies the norm recursion, while parity and coordinate phases supply the rest.
\writeup Moving the creation operator across the inner product gives
$\|g_n\|_2^2=\langle g_{n-1},Ag_n\rangle=4\pi n\|g_{n-1}\|_2^2$.
Since $\|g_0\|_2^2=1/\sqrt2$, induction yields $(4\pi)^nn!/\sqrt2$. Direct application gives
$g_1=4\pi xe^{-\pi x^2}$ and $g_2=(16\pi^2x^2-4\pi)e^{-\pi x^2}$.
The differentiation formula gives $\mathcal F(xe^{-\pi x^2})=\tfrac i{2\pi}(e^{-\pi\xi^2})'=-i\xi e^{-\pi\xi^2}$, proving the assertion for $g_1$.
The operator $A^\dagger$ switches even and odd functions. Starting with even $g_0$ proves the parity formula. Multiplying each Hermite coordinate by $e^{-in\alpha}$ preserves its squared modulus, so $\mathcal F_\alpha$ is unitary with inverse $\mathcal F_{-\alpha}$. Coordinate multiplication gives the group law. At $\alpha=\pi/2$ the phases are $(-i)^n$, agreeing with $\mathcal F$ on a complete basis; at $\alpha=\pi$ they are $(-1)^n$, agreeing with reflection by parity.
\end{solution}

\begin{solution}{ex:hysharp}
\idea Compute one norm ratio carefully, then distinguish that ratio from a supremum over all functions.
\writeup Gaussian integration gives $\|\gamma\|_p=p^{-1/(2p)}$, while self-duality gives $\|\hat\gamma\|_{p'}=p'^{-1/(2p')}$. Put $a=1/p\in(1/2,1)$. The logarithm of their ratio is
$\tfrac12((1-a)\log(1-a)-a\log a)$.
The function $E(a)=a\log a-(1-a)\log(1-a)$ vanishes at $1/2$ and $1$, with $E''(a)=1/a-1/(1-a)<0$ inside. Strict concavity gives $E(a)>0$, proving the ratio is less than one.
This calculation gives a lower bound for the Fourier operator norm, and shows only that the norm-one estimate has slack on this input. It cannot alone give an improved upper bound for every input. Beckner's quoted theorem supplies precisely that global bound: the operator norm is this ratio, $\bigl(p^{1/p}/p'^{1/p'}\bigr)^{1/2}$ in one dimension, raised to the power $d$ in dimension $d$. Its strictness therefore improves the constant, but the global assertion is an additional theorem, not a consequence of the single Gaussian calculation.
\end{solution}

\begin{solution}{ex:hexlattice}
\idea Inverse transpose gives the dual basis; symmetry then reduces the dual theta function to the original one.
\writeup The basis matrix and inverse transpose are
\[
A=\begin{pmatrix}1&1/2\\0&\sqrt3/2\end{pmatrix},\qquad
A^{-T}=\begin{pmatrix}1&0\\-1/\sqrt3&2/\sqrt3\end{pmatrix}.
\]
Thus the covolume is $\sqrt3/2$ and the dual has basis $w_1=(1,-1/\sqrt3)$, $w_2=(0,2/\sqrt3)$. Rotation $R(x,y)=(-y,x)$ followed by scale $2/\sqrt3$ sends the original basis to $w_2,-w_1$, another basis of that dual. Consequently $\theta_{\Lambda^*}(s)=\theta_\Lambda(4s/3)$, and the theta transformation gives $\theta_\Lambda(t)=2(\sqrt3t)^{-1}\theta_\Lambda(4/(3t))$.
At $t_0=2/\sqrt3$, the argument returns to $t_0$ and the prefactor is one, so direct substitution is a tautology. The sums and their derivatives converge locally uniformly for $t>0$, by Gaussian decay. Differentiating the identity at $t_0$ gives the additional relation $t_0\theta_\Lambda'(t_0)=-\theta_\Lambda(t_0)/2$.
\end{solution}

\begin{solution}{ex:v8c10a}
\idea Gaussian scaling controls both width and amplitude; its differentiated integral controls the second moments.
\writeup The transform is $\hat f_a(\xi)=a^{-1/2}e^{-\pi\xi^2/a}$ and $\|f_a\|_2=(2a)^{-1/4}$. Differentiating $\int e^{-2\pi ax^2}\dd x=(2a)^{-1/2}$ under the integral gives
$\int x^2e^{-2\pi ax^2}\dd x=(2a)^{-1/2}/(4\pi a)$.
The derivative is justified on any compact positive interval of parameters by an integrable Gaussian times $x^2$. Hence
\[
\|xf_a\|_2=\frac{\|f_a\|_2}{2\sqrt{\pi a}},\qquad
\|\xi\hat f_a\|_2=\frac{\sqrt a\,\|f_a\|_2}{2\sqrt\pi}.
\]
For the second identity apply the same Gaussian second-moment calculation to $a^{-1/2}e^{-\pi\xi^2/a}$, or combine it with Plancherel. Their product is $\|f_a\|_2^2/(4\pi)$, the asserted uncertainty equality.
\end{solution}

\begin{solution}{ex:v8c10b}
\idea At the specified bandwidth, the two summands are orthonormal sampling vectors.
\writeup At integers, $f(0)=1$, $f(1)=2$, and all other samples vanish. The inverse transform of $e^{-2\pi in\xi}\mathbf1_{[-1/2,1/2]}$ is $\operatorname{sinc}(x-n)$, so
$\hat f(\xi)=(1+2e^{-2\pi i\xi})\mathbf1_{[-1/2,1/2]}(\xi)$.
The two sinc terms are orthonormal, giving $\|f\|_2^2=1+4=5$. Since $\operatorname{sinc}(1/2)=\operatorname{sinc}(-1/2)=2/\pi$, their value at $1/2$ is $6/\pi$.
For general bandwidth, the coefficient of the normalised basis vector $\varphi_n$ is $g(n/(2W))/\sqrt{2W}$, as proved in the sampling theorem. Parseval gives the stated sample-energy identity, including convergence of its sum.
\end{solution}

\begin{solution}{ex:v8c10c}
\idea Atomic measures turn the transform into a finite exponential sum, with no decay requirement.
\writeup We have $\hat\mu(\xi)=(e^{2\pi i\xi}+e^{-2\pi i\xi})/2=\cos(2\pi\xi)$. In a sum of $n$ independent signs, choosing $k$ negative signs gives value $n-2k$ with multiplicity $\binom nk$. Therefore
\[
\mu^{*n}=2^{-n}\sum_{k=0}^n\binom nk\delta_{n-2k},\qquad
\widehat{\mu^{*n}}(\xi)=\cos(2\pi\xi)^n.
\]
At integer frequencies the original transform equals one, so it does not decay. This measure has atoms and no Lebesgue $L^1$ density. Riemann--Lebesgue asserted decay for transforms of $L^1$ functions, not for every finite measure. In fact its failure of decay also rules out the existence of any such density by uniqueness.
\end{solution}

\begin{solution}{ex:v8c10d}
\idea Separate coordinates first; rotation and scaling then identify the dual rectangular lattice.
\writeup The diagonal basis matrix has determinant $6$, and its inverse transpose gives $\Lambda^*=(1/2)\Z\times(1/3)\Z$. The two theta functions separate into $\theta_\Lambda(t)=\theta(4t)\theta(9t)$ and $\theta_{\Lambda^*}(1/t)=\theta(1/(4t))\theta(1/(9t))$. The lattice theorem gives exactly the displayed product identity. The leading term is $1/(6t)$ as $t\downarrow0$.
Let $R$ rotate by $90$ degrees. The lattice $(1/6)R\Lambda$ is $(1/2)\Z\times(1/3)\Z$, because the rotated basis vectors are $(0,1/3)$ and $(-1/2,0)$. Thus $\theta_{\Lambda^*}(s)=\theta_\Lambda(s/36)$, yielding $\theta_\Lambda(t)=(6t)^{-1}\theta_\Lambda(1/(36t))$. Its positive fixed point solves $t^2=1/36$, namely $t=1/6$, where the prefactor is also one.
\end{solution}

\section*{Sources and further reading}

\begin{itemize}
\item J.~Fourier, \emph{Th\'eorie analytique de la chaleur}, Didot, 1822; T.~W.~K\"orner, \emph{Fourier Analysis}, Cambridge, 1988, for historical context and worked applications.
\item M.~Plancherel, ``Contribution \`a l'\'etude de la repr\'esentation d'une fonction arbitraire par des int\'egrales d\'efinies'', \emph{Rend. Circ. Mat. Palermo} 30 (1910), 289--335.
\item M.~Riesz, ``Sur les maxima des formes bilin\'eaires et sur les fonctionnelles lin\'eaires'', \emph{Acta Math.} 49 (1927), 465--497; G.~O.~Thorin, ``An extension of a convexity theorem due to M.~Riesz'',
\emph{Kungl. Fysiogr. S\"allsk. i Lund F\"orh.} 8 (1938), 166--170; J.~E.~Littlewood, \emph{Littlewood's Miscellany}, ed. B.~Bollob\'as, Cambridge, 1986.
\item W.~Beckner, ``Inequalities in Fourier analysis'', \emph{Ann. of Math.} 102 (1975), 159--182.
\item E.~M.~Stein and G.~Weiss, \emph{Introduction to Fourier Analysis on Euclidean Spaces}, Princeton, 1971; E.~M.~Stein and R.~Shakarchi, \emph{Fourier Analysis}, Princeton, 2003.
\item R.~E.~A.~C.~Paley and N.~Wiener, \emph{Fourier Transforms in the Complex Domain}, AMS, 1934; C.~E.~Shannon, ``Communication in the presence of noise'', \emph{Proc. IRE} 37 (1949), 10--21.
\item M.~Viazovska, ``The sphere packing problem in dimension 8'', \emph{Ann. of Math.} 185 (2017), 991--1015.
\item I.~Grattan-Guinness, \emph{Joseph Fourier, 1768--1830}, MIT Press, 1972, reproduces the 1807 memoir and discusses the committee's response; J.~Herivel, \emph{Joseph Fourier: The Man and the Physicist}, Oxford, 1975. The cited \emph{Miscellany} records Littlewood's assessment of Thorin's argument.
\end{itemize}
